%QUESTÃO - Cód.: TF1p02q09

\question [\questaopt] Enquanto o professor escreve na lousa:
\begin{parts}

\part o giz está em repouso ou em movimento em relação à lousa?
\part a lousa está em repouso ou em movimento em relação ao chão?
\part a lousa está em repouso ou em movimento em relação ao giz?\\
 \\
\end{parts}


%--------------------------SOLUÇÃO---------------------------------------

	\begin{solution}
		
	(a) Enquanto o professor está escrevendo, o giz muda de posição em
	relação à lousa, estando, portanto, em movimento em relação a
	ela.\\
	(b) A lousa não muda de posição em relação ao chão, estando, por-
	tanto, em repouso em relação a ele.\\
	(c) Os conceitos de movimento e de repouso são simétricos, isto é, se
	um corpo está em movimento (ou repouso) em relação a outro,
	este também está em movimento (ou repouso) em relação ao primeiro. Assim, a lousa está em movimento em relação ao giz. De
	fato, se houver um inseto pousado no giz, por exemplo, o inseto
	verá a lousa passando por ele.
		
	\end{solution}

